\documentclass[10pt,a4paper]{amsart}
\usepackage{amsmath,amssymb,amsthm,graphicx}
% For commutative diagrams we use Xy-pic
\usepackage[all]{xy}

% We use 2cell for 2-commutative diagrams.
\xyoption{2cell}
% Theorem environments.
%
\theoremstyle{plain}
\newtheorem{theorem}[subsection]{Theorem}
\newtheorem{proposition}[subsection]{Proposition}
\newtheorem{lemma}[subsection]{Lemma}

\theoremstyle{definition}
\newtheorem{definition}[subsection]{Definition}
\newtheorem{example}[subsection]{Example}
\newtheorem{exercise}[subsection]{Exercise}
\newtheorem{situation}[subsection]{Situation}

\theoremstyle{remark}
\newtheorem{remark}[subsection]{Remark}
\newtheorem{remarks}[subsection]{Remarks}

\numberwithin{equation}{subsection}

% Macros
%
\def\lim{\mathop{\mathrm{lim}}\nolimits}
\def\colim{\mathop{\mathrm{colim}}\nolimits}
\def\Spec{\mathop{\mathrm{Spec}}}
\def\Hom{\mathop{\mathrm{Hom}}\nolimits}
\def\Ext{\mathop{\mathrm{Ext}}\nolimits}
\def\SheafHom{\mathop{\mathcal{H}\!\mathit{om}}\nolimits}
\def\SheafExt{\mathop{\mathcal{E}\!\mathit{xt}}\nolimits}
\def\Sch{\mathit{Sch}}
\def\Mor{\mathop{\mathrm{Mor}}\nolimits}
\def\Ob{\mathop{\mathrm{Ob}}\nolimits}
\def\Sh{\mathop{\mathit{Sh}}\nolimits}
\def\NL{\mathop{N\!L}\nolimits}
\def\CH{\mathop{\mathrm{CH}}\nolimits}
\def\proetale{{pro\text{-}\acute{e}tale}}
\def\etale{{\acute{e}tale}}
\def\QCoh{\mathit{QCoh}}
\def\Ker{\mathop{\mathrm{Ker}}}
\def\Im{\mathop{\mathrm{Im}}}
\def\Coker{\mathop{\mathrm{Coker}}}
\def\Coim{\mathop{\mathrm{Coim}}}

% Boxtimes
%
\DeclareMathSymbol{\boxtimes}{\mathbin}{AMSa}{"02}

%
% Macros for moduli stacks/spaces
%
\def\QCohstack{\mathcal{QC}\!\mathit{oh}}
\def\Cohstack{\mathcal{C}\!\mathit{oh}}
\def\Spacesstack{\mathcal{S}\!\mathit{paces}}
\def\Quotfunctor{\mathrm{Quot}}
\def\Hilbfunctor{\mathrm{Hilb}}
\def\Curvesstack{\mathcal{C}\!\mathit{urves}}
\def\Polarizedstack{\mathcal{P}\!\mathit{olarized}}
\def\Complexesstack{\mathcal{C}\!\mathit{omplexes}}
% \Pic is the operator that assigns to X its picard group, usage \Pic(X)
% \Picardstack_{X/B} denotes the Picard stack of X over B
% \Picardfunctor_{X/B} denotes the Picard functor of X over B
\def\Pic{\mathop{\mathrm{Pic}}\nolimits}
\def\Picardstack{\mathcal{P}\!\mathit{ic}}
\def\Picardfunctor{\mathrm{Pic}}
\def\Deformationcategory{\mathcal{D}\!\mathit{ef}}
\newenvironment{reference}{}{}
\newenvironment{slogan}{}{}
\newenvironment{history}{}{}
\usepackage{cleveref}
\crefname{originalsection}{original section}{original sections}
\Crefname{originalsection}{Original section}{Original sections}
\makeatletter
\newlabel{categories-remark-cones-and-cocones}{{0G2U}{source}{}{sourcepointer-1}{}}
\newlabel{categories-remark-cones-and-cocones@cref}{{[originalsection][0][]0G2U}{[1][0][]source}{}{}{}}
\makeatother
\crefname{originallemma}{original lemma}{original lemmas}
\Crefname{originallemma}{Original lemma}{Original lemmas}
\makeatletter
\newlabel{spaces-properties-lemma-cover-by-union-affines}{{03FX}{source}{}{sourcepointer-2}{}}
\newlabel{spaces-properties-lemma-cover-by-union-affines@cref}{{[originallemma][0][]03FX}{[1][0][]source}{}{}{}}
\makeatother
\crefname{originallemma}{original lemma}{original lemmas}
\Crefname{originallemma}{Original lemma}{Original lemmas}
\makeatletter
\newlabel{lemma-pushout-fpqc-local}{{0GFQ}{source}{}{sourcepointer-3}{}}
\newlabel{lemma-pushout-fpqc-local@cref}{{[originallemma][0][]0GFQ}{[1][0][]source}{}{}{}}
\makeatother
\crefname{originaldefinition}{original definition}{original definitions}
\Crefname{originaldefinition}{Original definition}{Original definitions}
\makeatletter
\newlabel{spaces-descent-definition-descent-datum}{{0ADG}{source}{}{sourcepointer-4}{}}
\newlabel{spaces-descent-definition-descent-datum@cref}{{[originaldefinition][0][]0ADG}{[1][0][]source}{}{}{}}
\makeatother
\crefname{originaldefinition}{original definition}{original definitions}
\Crefname{originaldefinition}{Original definition}{Original definitions}
\makeatletter
\newlabel{spaces-descent-definition-effective}{{0ADQ}{source}{}{sourcepointer-5}{}}
\newlabel{spaces-descent-definition-effective@cref}{{[originaldefinition][0][]0ADQ}{[1][0][]source}{}{}{}}
\makeatother
\crefname{originalsection}{original section}{original sections}
\Crefname{originalsection}{Original section}{Original sections}
\makeatletter
\newlabel{spaces-properties-section-etale-site}{{03EB}{source}{}{sourcepointer-6}{}}
\newlabel{spaces-properties-section-etale-site@cref}{{[originalsection][0][]03EB}{[1][0][]source}{}{}{}}
\makeatother
\crefname{originallemma}{original lemma}{original lemmas}
\Crefname{originallemma}{Original lemma}{Original lemmas}
\makeatletter
\newlabel{lemma-glue-etale-sheaf-proper-surjective}{{0GFW}{source}{}{sourcepointer-7}{}}
\newlabel{lemma-glue-etale-sheaf-proper-surjective@cref}{{[originallemma][0][]0GFW}{[1][0][]source}{}{}{}}
\makeatother
\crefname{originallemma}{original lemma}{original lemmas}
\Crefname{originallemma}{Original lemma}{Original lemmas}
\makeatletter
\newlabel{spaces-properties-lemma-sheaf-gives-space}{{0GF6}{source}{}{sourcepointer-8}{}}
\newlabel{spaces-properties-lemma-sheaf-gives-space@cref}{{[originallemma][0][]0GF6}{[1][0][]source}{}{}{}}
\makeatother
\crefname{originallemma}{original lemma}{original lemmas}
\Crefname{originallemma}{Original lemma}{Original lemmas}
\makeatletter
\newlabel{spaces-morphisms-lemma-base-change-proper}{{04WP}{source}{}{sourcepointer-9}{}}
\newlabel{spaces-morphisms-lemma-base-change-proper@cref}{{[originallemma][0][]04WP}{[1][0][]source}{}{}{}}
\makeatother
\crefname{originallemma}{original lemma}{original lemmas}
\Crefname{originallemma}{Original lemma}{Original lemmas}
\makeatletter
\newlabel{spaces-morphisms-lemma-composition-proper}{{04XY}{source}{}{sourcepointer-10}{}}
\newlabel{spaces-morphisms-lemma-composition-proper@cref}{{[originallemma][0][]04XY}{[1][0][]source}{}{}{}}
\makeatother
\crefname{originallemma}{original lemma}{original lemmas}
\Crefname{originallemma}{Original lemma}{Original lemmas}
\makeatletter
\newlabel{spaces-cohomology-lemma-proper-pushforward-coherent}{{08AR}{source}{}{sourcepointer-11}{}}
\newlabel{spaces-cohomology-lemma-proper-pushforward-coherent@cref}{{[originallemma][0][]08AR}{[1][0][]source}{}{}{}}
\makeatother
\crefname{originaldefinition}{original definition}{original definitions}
\Crefname{originaldefinition}{Original definition}{Original definitions}
\makeatletter
\newlabel{spaces-morphisms-definition-relative-spec}{{081W}{source}{}{sourcepointer-12}{}}
\newlabel{spaces-morphisms-definition-relative-spec@cref}{{[originaldefinition][0][]081W}{[1][0][]source}{}{}{}}
\makeatother
\crefname{originalsection}{original section}{original sections}
\Crefname{originalsection}{Original section}{Original sections}
\makeatletter
\newlabel{spaces-morphisms-remark-factorization-quasi-compact-quasi-separated}{{081X}{source}{}{sourcepointer-13}{}}
\newlabel{spaces-morphisms-remark-factorization-quasi-compact-quasi-separated@cref}{{[originalsection][0][]081X}{[1][0][]source}{}{}{}}
\makeatother
\crefname{originallemma}{original lemma}{original lemmas}
\Crefname{originallemma}{Original lemma}{Original lemmas}
\makeatletter
\newlabel{spaces-morphisms-lemma-universally-closed-permanence}{{04NX}{source}{}{sourcepointer-14}{}}
\newlabel{spaces-morphisms-lemma-universally-closed-permanence@cref}{{[originallemma][0][]04NX}{[1][0][]source}{}{}{}}
\makeatother
\crefname{originaltheorem}{original theorem}{original theorems}
\Crefname{originaltheorem}{Original theorem}{Original theorems}
\makeatletter
\newlabel{spaces-more-morphisms-theorem-stein-factorization-Noetherian}{{0A1B}{source}{}{sourcepointer-15}{}}
\newlabel{spaces-more-morphisms-theorem-stein-factorization-Noetherian@cref}{{[originaltheorem][0][]0A1B}{[1][0][]source}{}{}{}}
\makeatother
\crefname{originallemma}{original lemma}{original lemmas}
\Crefname{originallemma}{Original lemma}{Original lemmas}
\makeatletter
\newlabel{algebra-lemma-integral-overring-surjective}{{00GQ}{source}{}{sourcepointer-16}{}}
\newlabel{algebra-lemma-integral-overring-surjective@cref}{{[originallemma][0][]00GQ}{[1][0][]source}{}{}{}}
\makeatother
\crefname{originaldefinition}{original definition}{original definitions}
\Crefname{originaldefinition}{Original definition}{Original definitions}
\makeatletter
\newlabel{spaces-morphisms-definition-universally-injective}{{03MV}{source}{}{sourcepointer-17}{}}
\newlabel{spaces-morphisms-definition-universally-injective@cref}{{[originaldefinition][0][]03MV}{[1][0][]source}{}{}{}}
\makeatother
\crefname{originallemma}{original lemma}{original lemmas}
\Crefname{originallemma}{Original lemma}{Original lemmas}
\makeatletter
\newlabel{spaces-morphisms-lemma-universally-injective}{{040X}{source}{}{sourcepointer-18}{}}
\newlabel{spaces-morphisms-lemma-universally-injective@cref}{{[originallemma][0][]040X}{[1][0][]source}{}{}{}}
\makeatother
\crefname{originallemma}{original lemma}{original lemmas}
\Crefname{originallemma}{Original lemma}{Original lemmas}
\makeatletter
\newlabel{spaces-morphisms-lemma-base-change-surjective}{{03MH}{source}{}{sourcepointer-19}{}}
\newlabel{spaces-morphisms-lemma-base-change-surjective@cref}{{[originallemma][0][]03MH}{[1][0][]source}{}{}{}}
\makeatother
\crefname{originallemma}{original lemma}{original lemmas}
\Crefname{originallemma}{Original lemma}{Original lemmas}
\makeatletter
\newlabel{spaces-morphisms-lemma-composition-surjective}{{03MG}{source}{}{sourcepointer-20}{}}
\newlabel{spaces-morphisms-lemma-composition-surjective@cref}{{[originallemma][0][]03MG}{[1][0][]source}{}{}{}}
\makeatother
\crefname{originallemma}{original lemma}{original lemmas}
\Crefname{originallemma}{Original lemma}{Original lemmas}
\makeatletter
\newlabel{spaces-morphisms-lemma-locally-finite-type-locally-noetherian}{{04ZK}{source}{}{sourcepointer-21}{}}
\newlabel{spaces-morphisms-lemma-locally-finite-type-locally-noetherian@cref}{{[originallemma][0][]04ZK}{[1][0][]source}{}{}{}}
\makeatother
\crefname{originallemma}{original lemma}{original lemmas}
\Crefname{originallemma}{Original lemma}{Original lemmas}
\makeatletter
\newlabel{spaces-cohomology-lemma-flat-base-change-cohomology}{{073K}{source}{}{sourcepointer-22}{}}
\newlabel{spaces-cohomology-lemma-flat-base-change-cohomology@cref}{{[originallemma][0][]073K}{[1][0][]source}{}{}{}}
\makeatother
\crefname{originallemma}{original lemma}{original lemmas}
\Crefname{originallemma}{Original lemma}{Original lemmas}
\makeatletter
\newlabel{spaces-properties-lemma-morphism-to-affine-scheme}{{05Z1}{source}{}{sourcepointer-23}{}}
\newlabel{spaces-properties-lemma-morphism-to-affine-scheme@cref}{{[originallemma][0][]05Z1}{[1][0][]source}{}{}{}}
\makeatother
\title{A proper algebraic-space morphism has a finite universally injective kernel-pair coequalizer}
\author{The Stacks Project Authors}
\date{stacks-project:spaces-pushouts.tex:local-snapshot-20261011; GNU FDL 1.2 or later}
\begin{document}
\maketitle
\noindent\textit{Editorial note.} Public280 concerns constructible etale sheaf presentations and public293 is Beck monadicity. Neither supplies the proper algebraic-space kernel-pair coequalizer or its finite radicial factorization.\par\medskip
\noindent\textit{Allowed standard prerequisites (editorial).}\begin{itemize}
\item Explicit prerequisite boundary. Locally Noetherian algebraic spaces, proper coherent pushforward, relative Spec, Stein factorization and flat base change. The local etale-check and proper-surjective descent support are retained; the coequalizer itself is not assumed.
\end{itemize}\medskip
\ifdefined\hypertarget\hypertarget{sourcepointer-1}{}\fi
\noindent\textit{Source reference.} Original section 0G2U: focused-seeds/AGP/sources/stacks-snapshot-20261011/tags.txt, source lines 20069--20069. Stable Stacks tag for an expressly allowed foundational or notational input; the counted construction remains the two concrete core obligations. The retained local core is not replaced by this pointer. See standard\_prerequisites.\par\medskip
\ifdefined\hypertarget\hypertarget{sourcepointer-2}{}\fi
\noindent\textit{Source reference.} Original lemma 03FX: focused-seeds/AGP/sources/stacks-snapshot-20261011/tags.txt, source lines 4508--4508. Stable Stacks tag for an expressly allowed foundational or notational input; the counted construction remains the two concrete core obligations. The retained local core is not replaced by this pointer. See standard\_prerequisites.\par\medskip
\ifdefined\hypertarget\hypertarget{sourcepointer-3}{}\fi
\noindent\textit{Source reference.} Original lemma 0GFQ: focused-seeds/AGP/sources/stacks-snapshot-20261011/tags.txt, source lines 20520--20520. Stable Stacks tag for an expressly allowed foundational or notational input; the counted construction remains the two concrete core obligations. The retained local core is not replaced by this pointer. See standard\_prerequisites.\par\medskip
\ifdefined\hypertarget\hypertarget{sourcepointer-4}{}\fi
\noindent\textit{Source reference.} Original definition 0ADG: focused-seeds/AGP/sources/stacks-snapshot-20261011/tags.txt, source lines 13089--13089. Stable Stacks tag for an expressly allowed foundational or notational input; the counted construction remains the two concrete core obligations. The retained local core is not replaced by this pointer. See standard\_prerequisites.\par\medskip
\ifdefined\hypertarget\hypertarget{sourcepointer-5}{}\fi
\noindent\textit{Source reference.} Original definition 0ADQ: focused-seeds/AGP/sources/stacks-snapshot-20261011/tags.txt, source lines 13098--13098. Stable Stacks tag for an expressly allowed foundational or notational input; the counted construction remains the two concrete core obligations. The retained local core is not replaced by this pointer. See standard\_prerequisites.\par\medskip
\ifdefined\hypertarget\hypertarget{sourcepointer-6}{}\fi
\noindent\textit{Source reference.} Original section 03EB: focused-seeds/AGP/sources/stacks-snapshot-20261011/tags.txt, source lines 4450--4450. Stable Stacks tag for an expressly allowed foundational or notational input; the counted construction remains the two concrete core obligations. The retained local core is not replaced by this pointer. See standard\_prerequisites.\par\medskip
\ifdefined\hypertarget\hypertarget{sourcepointer-7}{}\fi
\noindent\textit{Source reference.} Original lemma 0GFW: focused-seeds/AGP/sources/stacks-snapshot-20261011/tags.txt, source lines 20526--20526. Stable Stacks tag for an expressly allowed foundational or notational input; the counted construction remains the two concrete core obligations. The retained local core is not replaced by this pointer. See standard\_prerequisites.\par\medskip
\ifdefined\hypertarget\hypertarget{sourcepointer-8}{}\fi
\noindent\textit{Source reference.} Original lemma 0GF6: focused-seeds/AGP/sources/stacks-snapshot-20261011/tags.txt, source lines 20501--20501. Stable Stacks tag for an expressly allowed foundational or notational input; the counted construction remains the two concrete core obligations. The retained local core is not replaced by this pointer. See standard\_prerequisites.\par\medskip
\ifdefined\hypertarget\hypertarget{sourcepointer-9}{}\fi
\noindent\textit{Source reference.} Original lemma 04WP: focused-seeds/AGP/sources/stacks-snapshot-20261011/tags.txt, source lines 6367--6367. Stable Stacks tag for an expressly allowed foundational or notational input; the counted construction remains the two concrete core obligations. The retained local core is not replaced by this pointer. See standard\_prerequisites.\par\medskip
\ifdefined\hypertarget\hypertarget{sourcepointer-10}{}\fi
\noindent\textit{Source reference.} Original lemma 04XY: focused-seeds/AGP/sources/stacks-snapshot-20261011/tags.txt, source lines 6411--6411. Stable Stacks tag for an expressly allowed foundational or notational input; the counted construction remains the two concrete core obligations. The retained local core is not replaced by this pointer. See standard\_prerequisites.\par\medskip
\ifdefined\hypertarget\hypertarget{sourcepointer-11}{}\fi
\noindent\textit{Source reference.} Original lemma 08AR: focused-seeds/AGP/sources/stacks-snapshot-20261011/tags.txt, source lines 10541--10541. Stable Stacks tag for an expressly allowed foundational or notational input; the counted construction remains the two concrete core obligations. The retained local core is not replaced by this pointer. See standard\_prerequisites.\par\medskip
\ifdefined\hypertarget\hypertarget{sourcepointer-12}{}\fi
\noindent\textit{Source reference.} Original definition 081W: focused-seeds/AGP/sources/stacks-snapshot-20261011/tags.txt, source lines 10231--10231. Stable Stacks tag for an expressly allowed foundational or notational input; the counted construction remains the two concrete core obligations. The retained local core is not replaced by this pointer. See standard\_prerequisites.\par\medskip
\ifdefined\hypertarget\hypertarget{sourcepointer-13}{}\fi
\noindent\textit{Source reference.} Original section 081X: focused-seeds/AGP/sources/stacks-snapshot-20261011/tags.txt, source lines 10232--10232. Stable Stacks tag for an expressly allowed foundational or notational input; the counted construction remains the two concrete core obligations. The retained local core is not replaced by this pointer. See standard\_prerequisites.\par\medskip
\ifdefined\hypertarget\hypertarget{sourcepointer-14}{}\fi
\noindent\textit{Source reference.} Original lemma 04NX: focused-seeds/AGP/sources/stacks-snapshot-20261011/tags.txt, source lines 6095--6095. Stable Stacks tag for an expressly allowed foundational or notational input; the counted construction remains the two concrete core obligations. The retained local core is not replaced by this pointer. See standard\_prerequisites.\par\medskip
\ifdefined\hypertarget\hypertarget{sourcepointer-15}{}\fi
\noindent\textit{Source reference.} Original theorem 0A1B: focused-seeds/AGP/sources/stacks-snapshot-20261011/tags.txt, source lines 12664--12664. Stable Stacks tag for an expressly allowed foundational or notational input; the counted construction remains the two concrete core obligations. The retained local core is not replaced by this pointer. See standard\_prerequisites.\par\medskip
\ifdefined\hypertarget\hypertarget{sourcepointer-16}{}\fi
\noindent\textit{Source reference.} Original lemma 00GQ: focused-seeds/AGP/sources/stacks-snapshot-20261011/tags.txt, source lines 616--616. Stable Stacks tag for an expressly allowed foundational or notational input; the counted construction remains the two concrete core obligations. The retained local core is not replaced by this pointer. See standard\_prerequisites.\par\medskip
\ifdefined\hypertarget\hypertarget{sourcepointer-17}{}\fi
\noindent\textit{Source reference.} Original definition 03MV: focused-seeds/AGP/sources/stacks-snapshot-20261011/tags.txt, source lines 4763--4763. Stable Stacks tag for an expressly allowed foundational or notational input; the counted construction remains the two concrete core obligations. The retained local core is not replaced by this pointer. See standard\_prerequisites.\par\medskip
\ifdefined\hypertarget\hypertarget{sourcepointer-18}{}\fi
\noindent\textit{Source reference.} Original lemma 040X: focused-seeds/AGP/sources/stacks-snapshot-20261011/tags.txt, source lines 5272--5272. Stable Stacks tag for an expressly allowed foundational or notational input; the counted construction remains the two concrete core obligations. The retained local core is not replaced by this pointer. See standard\_prerequisites.\par\medskip
\ifdefined\hypertarget\hypertarget{sourcepointer-19}{}\fi
\noindent\textit{Source reference.} Original lemma 03MH: focused-seeds/AGP/sources/stacks-snapshot-20261011/tags.txt, source lines 4749--4749. Stable Stacks tag for an expressly allowed foundational or notational input; the counted construction remains the two concrete core obligations. The retained local core is not replaced by this pointer. See standard\_prerequisites.\par\medskip
\ifdefined\hypertarget\hypertarget{sourcepointer-20}{}\fi
\noindent\textit{Source reference.} Original lemma 03MG: focused-seeds/AGP/sources/stacks-snapshot-20261011/tags.txt, source lines 4748--4748. Stable Stacks tag for an expressly allowed foundational or notational input; the counted construction remains the two concrete core obligations. The retained local core is not replaced by this pointer. See standard\_prerequisites.\par\medskip
\ifdefined\hypertarget\hypertarget{sourcepointer-21}{}\fi
\noindent\textit{Source reference.} Original lemma 04ZK: focused-seeds/AGP/sources/stacks-snapshot-20261011/tags.txt, source lines 6468--6468. Stable Stacks tag for an expressly allowed foundational or notational input; the counted construction remains the two concrete core obligations. The retained local core is not replaced by this pointer. See standard\_prerequisites.\par\medskip
\ifdefined\hypertarget\hypertarget{sourcepointer-22}{}\fi
\noindent\textit{Source reference.} Original lemma 073K: focused-seeds/AGP/sources/stacks-snapshot-20261011/tags.txt, source lines 9064--9064. Stable Stacks tag for an expressly allowed foundational or notational input; the counted construction remains the two concrete core obligations. The retained local core is not replaced by this pointer. See standard\_prerequisites.\par\medskip
\ifdefined\hypertarget\hypertarget{sourcepointer-23}{}\fi
\noindent\textit{Source reference.} Original lemma 05Z1: focused-seeds/AGP/sources/stacks-snapshot-20261011/tags.txt, source lines 7675--7675. Stable Stacks tag for an expressly allowed foundational or notational input; the counted construction remains the two concrete core obligations. The retained local core is not replaced by this pointer. See standard\_prerequisites.\par\medskip
\begin{lemma}
\label{lemma-coequalizer}
Let $S$ be a scheme. Let
$$
g : Y \longrightarrow X
$$
be a morphism of algebraic spaces over $S$. Assume $X$ is locally Noetherian,
and $g$ is proper. Let $R = Y \times_X Y$ with projection morphisms
$t, s : R \to Y$. There exists a coequalizer $X'$ of $s, t : R \to Y$
in the category of algebraic spaces over $S$. Moreover
\begin{enumerate}
\item The morphism $X' \to X$ is finite.
\item The morphism $Y \to X'$ is proper.
\item The morphism $Y \to X'$ is surjective.
\item The morphism $X' \to X$ is universally injective.
\item If $g$ is surjective, the morphism $X' \to X$
is a universal homeomorphism.
\end{enumerate}
\end{lemma}

\begin{lemma}
\label{lemma-colimit-check-etale-locally}
Let $S$ be a scheme. Let $B$ be an algebraic space over $S$.
Let $\mathcal{I} \to (\Sch/S)_{fppf}$, $i \mapsto X_i$
be a diagram of algebraic spaces over $B$. Let $(X, X_i \to X)$
be a cocone for the diagram in the category of algebraic spaces over $B$
(Categories, Remark \ref{categories-remark-cones-and-cocones}).
Assume that
\begin{enumerate}
\item the base change functor
$X_{spaces, \etale} \to \lim X_{i, spaces, \etale}$,
sending $U$ to $U_i = X_i \times_X U$ is an equivalence,
\item given
\begin{enumerate}
\item $B'$ affine and \'etale over $B$,
\item $Z$ an affine scheme over $B'$,
\item $U \to X \times_B B'$ an \'etale morphism of algebraic spaces
with $U$ affine,
\item $f_i : U_i \to Z$ a cocone over $B'$ of the diagram
$i \mapsto U_i = U \times_X X_i$,
\end{enumerate}
there exists a unique morphism $f : U \to Z$ over $B'$
such that $f_i$ equals the composition $U_i \to U \to Z$.
\end{enumerate}
Then $X = \colim X_i$ in the category of all algebraic spaces over $B$.
\end{lemma}

\begin{proof}
In this paragraph we reduce to the case where $B$ is an affine scheme.
Let $B' \to B$ be an \'etale morphism of algebraic spaces.
Observe that conditions (1) and (2) are preserved if we replace
$B$, $X_i$, $X$ by $B'$, $X_i \times_B B'$, $X \times_B B'$.
Let $\{B_a \to B\}_{a \in A}$ be an \'etale covering with $B_a$
affine, see Properties of Spaces, Lemma
\ref{spaces-properties-lemma-cover-by-union-affines}.
For $a \in A$ denote $X_a$, $X_{a, i}$ the base changes of $X$ and the
diagram to $B_a$. For $a, b \in A$ denote
$X_{a, b}$ and $X_{a, b, i}$ the base changes of $X$ and the
diagram to $B_a \times_B B_b$.
By Lemma \ref{lemma-pushout-fpqc-local}
it suffices to prove that $X_a = \colim X_{a, i}$ and
$X_{a, b} = \colim X_{a, b, i}$.
This reduces us to the case where $B = B_a$ (an affine scheme) or
$B = B_a \times_B B_b$ (a separated scheme). Repeating the
argument once more, we conclude that we may assume $B$ is
an affine scheme (this uses that the intersection of affine
opens in a separated scheme is affine).

\medskip\noindent
Assume $B$ is an affine scheme. Let $Z$ be an algebraic space over $B$.
We have to show
$$
\Mor_B(X, Z) \longrightarrow \lim \Mor_B(X_i, Z)
$$
is a bijection.

\medskip\noindent
Proof of injectivity. Let $f, g : X \to Z$ be morphisms
such that the compositions $f_i, g_i : X_i \to Z$ are the same for all $i$.
Choose an affine scheme $Z'$ and an \'etale morphism $Z' \to Z$. By
Properties of Spaces, Lemma \ref{spaces-properties-lemma-cover-by-union-affines}
we know we can cover $Z$ by such affines. Set $U = X \times_{f, Z} Z'$
and $U' = X \times_{g, Z} Z'$ and denote $p : U \to X$ and $p' : U' \to X$
the projections. Since $f_i = g_i$ for all $i$, we see that
$$
U_i = X_i \times_{f_i, Z} Z' = X_i \times_{g_i, Z} Z' = U'_i
$$
compatible with transition morphisms. By (1) there is a unique isomorphism
$\epsilon : U \to U'$ as algebraic spaces over $X$, i.e., with
$p = p' \circ \epsilon$ which is compatible with
the displayed identifications. Choose an \'etale covering
$\{h_a : U_a \to U\}$ with $U_a$ affine.
By (2) we see that $f \circ p \circ h_a = g \circ p' \circ \epsilon \circ h_a =
g \circ p \circ h_a$. Since $\{h_a : U_a \to U\}$ is an \'etale covering
we conclude $f \circ p = g \circ p$.
Since the collection of morphisms $p : U \to X$ we obtain
in this manner is an \'etale covering, we conclude that $f = g$.

\medskip\noindent
Proof of surjectivity. Let $f_i : X_i \to Z$ be an element of the
right hand side of the displayed arrow in the first paragraph of the proof.
It suffices to find an \'etale covering $\{U_c \to X\}_{c \in C}$
such that the families $f_{c, i} \in \lim_i \Mor_B(X_i \times_X U_c, Z)$
come from morphisms $f_c : U_c \to Z$. Namely, by the uniqueness
proved above the morphisms $f_c$ will agree on $U_c \times_X U_b$
and hence will descend to give the desired morphism $f : X \to Z$.
To find our covering, we first choose an \'etale covering
$\{g_a : Z_a \to Z\}_{a \in A}$ where each $Z_a$ is affine. Then we let
$U_{a, i} = X_i \times_{f_i, Z} Z_a$. By (1) we find
$U_{a, i} = X_i \times_X U_a$ for some algebraic spaces $U_a$ \'etale
over $X$. Then we choose \'etale coverings
$\{U_{a, b} \to U_a\}_{b \in B_a}$
with $U_{a, b}$ affine and we consider the morphisms
$$
U_{a, b, i} = X_i \times_X U_{a, b} \to
X_i \times_X U_a = X_i \times_{f_i, Z} Z_a \to Z_a
$$
By (2) we obtain morphisms $f_{a, b} : U_{a, b} \to Z_a$ compatible with these
morphisms. Setting $C = \coprod_{a \in A} B_a$ and for $c \in C$
corresponding to $b \in B_a$ setting $U_c = U_{a, b}$ and
$f_c = g_a \circ f_{a, b} : U_c \to Z$ we conclude.
\end{proof}

\begin{lemma}
\label{lemma-descend-etale-proper-surjective}
Let $S$ be a scheme. Let $f : X \to Y$ be a proper surjective morphism
of algebraic spaces over $S$. Any descent datum $(U/X, \varphi)$ relative to $f$
(Descent on Spaces, Definition \ref{spaces-descent-definition-descent-datum})
with $U$ \'etale over $X$ is effective
(Descent on Spaces, Definition \ref{spaces-descent-definition-effective}).
More precisely, there exists an \'etale morphism $V \to Y$ of algebraic spaces
whose corresponding canonical descent datum is isomorphic to $(U/X, \varphi)$.
\end{lemma}

\begin{proof}
Recall that $U$ gives rise to a representable sheaf
$\mathcal{F} = h_U$ in $\Sh(X_{spaces, \etale}) = \Sh(X_\etale)$, see
Properties of Spaces, Section \ref{spaces-properties-section-etale-site}.
The descent datum on $U$ relative to $f$
exactly gives a descent datum $(\mathcal{F}, \varphi)$
for \'etale sheaves with respect to $\{X \to Y\}$.
By Lemma \ref{lemma-glue-etale-sheaf-proper-surjective}
this descent datum is effective.
Let $\mathcal{G}$ be the corresponding sheaf on $Y_\etale$.
By Properties of Spaces, Lemma \ref{spaces-properties-lemma-sheaf-gives-space}
we obtain an \'etale morphism $V \to Y$ of algebraic spaces
corresponding to $\mathcal{G}$; we omit the verification of
the set theoretic condition\footnote{It follows from the
fact that $\mathcal{F}$ satisfies the corresponding condition.}.
The given isomorphism $\mathcal{F} \to f_{small}^{-1}\mathcal{G}$
corresponds to an isomorphism $U \to V \times_Y X$ compatible
with the descent datum.
\end{proof}

\begin{proof}
Denote $h : R \to X$ denote the composition of either $s$ or $t$
with $g$. Then $h$ is proper by Morphisms of Spaces, Lemmas
\ref{spaces-morphisms-lemma-base-change-proper} and
\ref{spaces-morphisms-lemma-composition-proper}.
The sheaves
$$
g_*\mathcal{O}_Y
\quad\text{and}\quad
h_*\mathcal{O}_R
$$
are coherent $\mathcal{O}_X$-algebras by Cohomology of Spaces, Lemma
\ref{spaces-cohomology-lemma-proper-pushforward-coherent}.
The $X$-morphisms $s$, $t$ induce $\mathcal{O}_X$-algebra maps
$s^\sharp, t^\sharp$ from the first to the second.
Set
$$
\mathcal{A} = \text{Equalizer}\left(s^\sharp, t^\sharp :
g_*\mathcal{O}_Y \longrightarrow h_*\mathcal{O}_R\right)
$$
Then $\mathcal{A}$ is a coherent $\mathcal{O}_X$-algebra and we
can define
$$
X' = \underline{\Spec}_X(\mathcal{A})
$$
as in Morphisms of Spaces, Definition
\ref{spaces-morphisms-definition-relative-spec}.
By Morphisms of Spaces, Remark
\ref{spaces-morphisms-remark-factorization-quasi-compact-quasi-separated}
and functoriality of the $\underline{\Spec}$ construction
there is a factorization
$$
Y \longrightarrow X' \longrightarrow X
$$
and the morphism $g' : Y \to X'$ equalizes $s$ and $t$.

\medskip\noindent
Before we show that $X'$ is the coequalizer of $s$ and $t$, we show
that $Y \to X'$ and $X' \to X$ have the desired properties. Since $\mathcal{A}$
is a coherent $\mathcal{O}_X$-module it is clear that $X' \to X$ is a
finite morphism of algebraic spaces. This proves (1).
The morphism $Y \to X'$ is proper by Morphisms of Spaces, Lemma
\ref{spaces-morphisms-lemma-universally-closed-permanence}. This proves (2).
Denote $Y \to Y' \to X$ with
$Y' = \underline{\Spec}_X(g_*\mathcal{O}_Y)$ the Stein factorization of $g$,
see More on Morphisms of Spaces, Theorem
\ref{spaces-more-morphisms-theorem-stein-factorization-Noetherian}.
Of course we obtain morphisms $Y \to Y' \to X' \to X$ fitting
with the morphisms studied above.
Since $\mathcal{O}_{X'} \subset g_*\mathcal{O}_Y$ is a finite extension
we see that $Y' \to X'$ is finite and surjective.
Some details omitted; hint: use
Algebra, Lemma \ref{algebra-lemma-integral-overring-surjective}
and reduce to the affine case by \'etale localization.
Since $Y \to Y'$ is surjective (with geometrically connected fibres)
we conclude that $Y \to X'$ is surjective. This proves (3).
To show that $X' \to X$ is universally injective, we have to show
that $X' \to X' \times_X X'$ is surjective, see
Morphisms of Spaces, Definition
\ref{spaces-morphisms-definition-universally-injective} and
Lemma \ref{spaces-morphisms-lemma-universally-injective}.
Since $Y \to X'$ is surjective (see above) and since base changes
and compositions of surjective morphisms are surjective by
Morphisms of Spaces, Lemmas
\ref{spaces-morphisms-lemma-base-change-surjective} and
\ref{spaces-morphisms-lemma-composition-surjective}
we see that $Y \times_X Y \to X' \times_X X'$ is surjective.
However, since $Y \to X'$ equalizes $s$ and $t$, we see that
$Y \times_X Y \to X' \times_X X'$ factors through $X' \to X' \times_X X'$
and we conclude this latter map is surjective. This proves (4).
Finally, if $g$ is surjective, then since $g$ factors through $X' \to X$
we see that $X' \to X$ is surjective. Since a surjective, universally
injective, finite morphism is a universal homeomorphism (because it
is universally bijective and universally closed), this proves (5).

\medskip\noindent
In the rest of the proof we show that $Y \to X'$ is the coequalizer
of $s$ and $t$ in the category of algebraic spaces over $S$.
Observe that $X'$ is locally Noetherian
(Morphisms of Spaces, Lemma
\ref{spaces-morphisms-lemma-locally-finite-type-locally-noetherian}).
Moreover, observe that $Y \times_{X'} Y \to Y \times_X Y$ is an
isomorphism as $Y \to X'$ equalizes $s$ and $t$ (this is a categorical
statement). Hence in order to prove the statement that $Y \to X'$
is the coequalizer of $s$ and $t$, we may and do assume $X = X'$.
In other words, $\mathcal{O}_X$ is the equalizer of the maps
$s^\sharp, t^\sharp : g_*\mathcal{O}_Y \to h_*\mathcal{O}_R$.

\medskip\noindent
Let $X_1 \to X$ be a flat morphism of algebraic spaces over $S$ with $X_1$
locally Noetherian. Denote $g_1 : Y_1 \to X_1$, $h_1 : R_1 \to X_1$ and
$s_1, t_1 : R_1 \to Y_1$ the base changes of $g, h, s, t$ to $X_1$.
Of course $g_1$ is proper and $R_1 = Y_1 \times_{X_1} Y_1$.
Since we have flat base change for pushforward of quasi-coherent modules,
Cohomology of Spaces, Lemma
\ref{spaces-cohomology-lemma-flat-base-change-cohomology}, we see that
$\mathcal{O}_{X_1}$ is the equalizer of the maps
$s_1^\sharp, t_1^\sharp : g_{1, *}\mathcal{O}_{Y_1} \to
h_{1, *}\mathcal{O}_{R_1}$. Hence all the assumptions we have
are preserved by this base change.

\medskip\noindent
At this point we are going to check conditions (1) and (2) of
Lemma \ref{lemma-colimit-check-etale-locally}. Condition (1)
follows from Lemma \ref{lemma-descend-etale-proper-surjective}
and the fact that $g$ is proper and surjective (because $X = X'$).
To check condition (2), by the remarks on base change above,
we reduce to the statement discussed and proved in the next paragraph.

\medskip\noindent
Assume $S = \Spec(A)$ is an affine scheme, $X = X'$ is an affine scheme, and
$Z$ is an affine scheme over $S$. We have to show that
$$
\Mor_S(X, Z) \longrightarrow
\text{Equalizer}(s, t : \Mor_S(Y, Z) \to \Mor_S(R, Z))
$$
is bijective. However, this is clear from the fact that $X = X'$
which implies $\mathcal{O}_X$ is the equalizer of the maps
$s^\sharp, t^\sharp : g_*\mathcal{O}_Y \to h_*\mathcal{O}_R$
which in turn implies
$$
\Gamma(X, \mathcal{O}_X) =
\text{Equalizer}\left(
s^\sharp, t^\sharp : \Gamma(Y, \mathcal{O}_Y) \to
\Gamma(R, \mathcal{O}_R)
\right)
$$
Namely, we have
$$
\Mor_S(X, Z) = \Hom_A(\Gamma(Z, \mathcal{O}_Z), \Gamma(X, \mathcal{O}_X))
$$
and similarly for $Y$ and $R$, see Properties of Spaces, Lemma
\ref{spaces-properties-lemma-morphism-to-affine-scheme}.
\end{proof}

\bibliographystyle{plain}
\bibliography{my}
\end{document}
